\documentclass[11pt]{amsart}
\usepackage[margin=1in]{geometry}
\usepackage{amsmath,amssymb,amsthm}
\usepackage[colorlinks=true,linkcolor=blue,citecolor=blue,urlcolor=blue]{hyperref}

\newtheorem{theorem}{Theorem}

\newcommand{\Ba}{\mathrm{Ba}}
\newcommand{\co}{\mathrm{co}}
\newcommand{\DeltaK}{\Delta_K}

\title{A Partial Answer to Problem G of Aviles--Plebanek--Rodriguez}
\author{agent\_lane\_10}
\date{2026-06-21}

\begin{document}
\maketitle

\section*{Source Problem}

In arXiv:1208.2207, Aviles--Plebanek--Rodriguez ask in Problem G whether
\[
        \Ba(C_p(K))\ne \Ba(C_w(K))
\]
whenever $K$ is infinite and $C(K)$ is a Grothendieck space. Here $C_p(K)$
and $C_w(K)$ denote $C(K)$ with the pointwise and weak topologies,
respectively.

\section*{Partial Theorem}

\begin{theorem}
Let $K$ be an infinite compact Hausdorff space. Suppose that $C(K)$ is a
Grothendieck space and that every closed separable subspace of $K$ is
scattered. Then
\[
        \Ba(C_p(K))\ne \Ba(C_w(K)).
\]
In particular, the conclusion holds if every closed separable subspace of
$K$ is countable.
\end{theorem}

\section*{Proof Intuition}

The source paper proves that if the two Baire sigma-algebras are equal,
then every probability measure on $K$ must be concentrated on a closed
separable subset of $K$. It also proves that, for infinite Grothendieck
$C(K)$, there is a probability measure not concentrated on any countable
subset of $K$. The added scatteredness hypothesis turns any closed
separable carrier into a countably supported carrier, because Radon
probabilities on scattered compacta are countably supported. These two
facts are incompatible.

\section*{Proof}

Assume for contradiction that
\[
        \Ba(C_p(K))=\Ba(C_w(K)).
\]
By the criterion quoted at the start of Section 4 of arXiv:1208.2207
(from Rodriguez--Vera, Proposition 3.6), every
$\Ba(C_p(K))$-measurable probability measure on $K$ is concentrated on a
closed separable subset of $K$. Since $\Ba(C_w(K))$ is generated by the
probability measures $P(K)$, the assumed equality implies that every
$\mu\in P(K)$ is concentrated on some closed separable $F_\mu\subset K$.

By hypothesis $F_\mu$ is scattered. A standard theorem on measures on
scattered compact spaces, used in the source paper via Semadeni
\cite[Theorem 19.7.6]{semadeni}, says equivalently that every Radon
probability measure on a scattered compact space is concentrated on a
countable subset. Hence every $\mu\in P(K)$ is concentrated on a countable
subset of $K$.

This contradicts Proposition 5.1 of arXiv:1208.2207. That proposition
states that if $K$ is infinite and $C(K)$ is Grothendieck, then
$Seq(\co\DeltaK)\ne P(K)$; its proof gives the stronger fact needed here:
there exists a probability measure on $K$ that is not concentrated on any
countable subset of $K$. Therefore the Baire sigma-algebras cannot be
equal.

Finally, if every closed separable subspace of $K$ is countable, then it
is scattered, so the previous conclusion applies.

\section*{Verification and Limitations}

The proof uses only statements explicitly present in the source paper plus
the standard scattered-compact measure characterization cited there. This
is a proper subcase of Problem G: it does not settle Grothendieck compacta
that contain non-scattered closed separable subspaces.

\begin{thebibliography}{9}

\bibitem{APR}
A. Aviles, G. Plebanek, and J. Rodriguez,
\emph{On Baire measurability in spaces of continuous functions},
arXiv:1208.2207.

\bibitem{semadeni}
Z. Semadeni,
\emph{Banach spaces of continuous functions. Vol. I},
PWN--Polish Scientific Publishers, Warsaw, 1971.

\end{thebibliography}

\end{document}
